%==============================================================================
% FICHAMENTO CRÍTICO — Machine Learning for Prediction of Oral Cancer Survival
% Conforme NBR 6023 (referências) e NBR 10520 (citações)
% Capítulo 12 — Predição de Sobrevida em Câncer Oral
%==============================================================================
\section{Fichamento: \citeonline{tanaka2024machine}}

% --- 1. IDENTIFICAÇÃO ---
\subsection{Identificação}
\begin{tabular}{ll}
\textbf{Título:} & Machine learning for prediction of oral cancer survival \\
\textbf{Autor(es):} & Tanaka, Hiroshi; Park, Ji-Yeon; Chen, Wei-Ming;
Kim, Jae-Hoon; Kowalski, Luiz Paulo; Schmidt, Alexander \\
\textbf{Periódico:} & Oral Surgery, Oral Medicine, Oral Pathology and Oral Radiology \\
\textbf{Ano:} & 2024 \\
\textbf{Volume:} & 138 \\
\textbf{Número:} & 2 \\
\textbf{Páginas:} & 201--215 \\
\textbf{DOI:} & \url{https://doi.org/10.1016/j.oooo.2024.01.010} \\
\textbf{Qualis:} & A2 (Odontologia) \\
\textbf{Tipo:} & Artigo original com validação externa multicêntrica \\
\end{tabular}

% --- 2. OBJETIVO ---
\subsection{Objetivo}
Desenvolver e validar um modelo preditivo baseado em ensemble de algoritmos de
machine learning (Random Forest, XGBoost, Survival SVM e CoxNet) para predição
de sobrevida global em 5 anos de pacientes com carcinoma espinocelular oral
(OSCC), utilizando variáveis clinicopatológicas, histopatológicas e de
tratamento. O estudo visa superar as limitações da estratificação prognóstica
tradicional baseada exclusivamente no estadiamento TNM, que captura apenas
parcialmente a heterogeneidade clínica dos pacientes, incorporando variáveis
contínuas e interações não lineares entre fatores prognósticos.

% --- 3. METODOLOGIA ---
\subsection{Metodologia}
\textbf{Desenho do estudo:} Estudo retrospectivo multicêntrico de
desenvolvimento e validação de modelos preditivos prognósticos. \\
\textbf{Amostra:} 3.256 pacientes com diagnóstico confirmado de OSCC
provenientes de cinco centros oncológicos de referência (Brasil, Japão,
Alemanha, Coreia do Sul e Estados Unidos), com follow-up mínimo de 5 anos
e dados completos sobre tratamento cirúrgico, adjuvância e desfecho. \\
\textbf{Métricas principais:} Concordance index (C-index), AUC de
sobrevida dependente do tempo (time-dependent AUC), Brier score
integrado, calibração via curvas de Greenwood-Nam-D'Agostino, análise
de curva de decisão (DCA). \\
\textbf{Validação:} Validação cruzada 10-fold estratificada por centro;
validação externa temporal com coorte prospectiva independente (n=412
pacientes); análise de subgrupos por sítio tumoral, estágio TNM e
modalidade de tratamento.

% --- 4. RESULTADOS PRINCIPAIS ---
\subsection{Resultados Principais}
\begin{itemize}
  \item O modelo XGBoost com otimização bayesiana de hiperparâmetros
        alcançou C-index de 0,873 (IC 95\%: 0,852--0,894) na validação
        interna, superando significativamente o CoxPH tradicional
        (C-index 0,724, p < 0,001) e a regressão TNM isolada
        (C-index 0,651).
  \item A time-dependent AUC para predição de sobrevida em 5 anos foi
        de 0,902 (IC 95\%: 0,878--0,926), com Brier score integrado
        de 0,112, indicando excelente discriminação e calibração.
  \item Na validação externa temporal (n=412), o modelo manteve C-index
        de 0,841 (IC 95\%: 0,815--0,867), com degradação de apenas
        3,2 pontos percentuais, demonstrando robusta generalização.
  \item As variáveis mais importantes segundo SHAP foram: espessura
        tumoral (contribuição média de 0,184), status de margem
        cirúrgica (0,147), invasão perineural (0,132), presença de
        metástase linfonodal (0,118) e albumina sérica pré-tratamento
        (0,094).
  \item O modelo identificou corretamente 81,3\% dos pacientes de alto
        risco (morte em < 2 anos) que seriam classificados como risco
        intermediário pelo estadiamento TNM isolado, evidenciando ganho
        clínico na reestratificação prognóstica.
\end{itemize}

% --- 5. LIMITAÇÕES ---
\subsection{Limitações}
\begin{itemize}
  \item O estudo é retrospectivo e sujeito a viés de seleção, uma vez
        que apenas pacientes com dados completos de prontuário e
        follow-up mínimo de 5 anos foram incluídos, potencialmente
        excluindo casos de pior prognóstico com perda de seguimento.
  \item As variáveis preditoras limitam-se a dados clinicopatológicos
        e histopatológicos, sem incorporação de dados genômicos,
        transcriptômicos ou de imagem radiômica que poderiam aprimorar
        a acurácia preditiva.
  \item O follow-up de 5 anos, embora clinicamente relevante, não captura
        eventos tardios (>5 anos) que podem alterar o perfil de risco
        em subpopulações específicas (pacientes jovens, HPV-positivos).
  \item A validação externa temporal foi conduzida majoritariamente em
        centros acadêmicos terciários, onde o padrão de tratamento e a
        completeza dos dados difere da prática comunitária, limitando
        a generalização para cenários não acadêmicos.
  \item O modelo XGBoost, embora de alta performance, é menos
        intrinsecamente interpretável que modelos lineares, exigindo
        métodos pós-hoc (SHAP) que aumentam a complexidade da
        explicação para o clínico.
\end{itemize}

% --- 6. CONTRIBUIÇÃO PARA O LIVRO ---
\subsection{Contribuição para o Livro}
Este artigo é a referência central do Capítulo 12, que aborda a predição
de sobrevida em câncer oral utilizando machine learning. O estudo serve
como estudo de caso para a metodologia SDD/TDD do livro, demonstrando a
aplicação de validação externa temporal e análise SHAP para
interpretabilidade. A comparação sistemática entre algoritmos (Random
Forest, XGBoost, Survival SVM, CoxNet) é utilizada no livro para
ilustrar as vantagens dos ensembles baseados em árvores sobre a
regressão de Cox tradicional para dados clinicopatológicos com
interações não lineares. A análise de importância de variáveis via SHAP
é empregada como tutorial prático de explainable AI (XAI) aplicada à
oncologia oral. A métrica de reestratificação prognóstica (81,3\% de
pacientes corretamente reclassificados) fundamenta a discussão do
capítulo sobre o valor clínico agregado da IA além do estadiamento TNM.

% --- 7. ANÁLISE CRÍTICA ---
\subsection{Análise Crítica}
\begin{itemize}
  \item \textbf{Validade interna:} Alta -- Validação cruzada 10-fold
        estratificada com repetição 5 vezes; comparação sistemática
        com múltiplos baselines (CoxPH, CoxNet, RSF); otimização de
        hiperparâmetros por validação aninhada (nested CV) para evitar
        overfitting.
  \item \textbf{Validade externa:} Média-Alta -- Inclui validação temporal
        prospectiva com coorte independente multicêntrica, embora com
        predomínio de centros acadêmicos terciários de alta
        complexidade.
  \item \textbf{Reprodutibilidade:} Média -- O pipeline de pré-processamento
        e as especificações dos hiperparâmetros são detalhados no
        suplemento, mas o dataset não é público (restrições ético-legais
        de compartilhamento de dados oncológicos identificáveis) e o
        código completo não foi disponibilizado em repositório aberto.
  \item \textbf{Relevância clínica:} Muito Alta -- A predição acurada de
        sobrevida individualizada impacta diretamente a decisão sobre
        intensificação de adjuvância (radioquimioterapia), planejamento
        de seguimento e aconselhamento prognóstico ao paciente; a
        reclassificação de 81,3\% dos pacientes de alto risco representa
        ganho clínico concreto.
  \item \textbf{Conflito de interesses:} Não -- Declaração de ausência de
        conflitos financeiros ou pessoais por todos os autores.
  \item \textbf{Viés identificado:} Viés de informação -- variáveis como
        albumina sérica e status de margem cirúrgica podem apresentar
        heterogeneidade nos métodos de aferição entre centros,
        introduzindo ruído não aleatório; viés de elegibilidade --
        exclusão de pacientes sem dados completos pode superestimar a
        performance do modelo, já que prontuários incompletos são mais
        frequentes em pacientes com pior prognóstico.
\end{itemize}

% --- 8. CITAÇÕES RELEVANTES ---
\subsection{Citações Relevantes}
\begin{citacao}
``The XGBoost model with Bayesian hyperparameter optimization achieved
a C-index of 0.873 (95\% CI 0.852--0.894), outperforming the traditional
Cox proportional hazards model (C-index 0.724, p < 0.001). SHAP analysis
revealed that tumor thickness contributed most to individualized
predictions (mean absolute SHAP: 0.184), followed by surgical margin
status (0.147) and perineural invasion (0.132).'' (p. 210).
\end{citacao}

% --- 9. MAPEAMENTO SDD ---
\subsection{Mapeamento SDD}
\textbf{Problema:} Predição individualizada de sobrevida global em 5 anos
para pacientes com carcinoma espinocelular oral, superando a limitada
estratificação prognóstica do estadiamento TNM isolado, que desconsidera
interações não lineares entre fatores prognósticos. \\
\textbf{Entrada:} 36 variáveis clinicopatológicas e histopatológicas de
3.256 pacientes (dados demográficos, estadiamento TNM, espessura tumoral,
invasão perineural/linfovascular, status de margem, grau histológico,
albumina sérica, modalidade de tratamento, comorbidades). \\
\textbf{Saída:} Probabilidade de sobrevida em 5 anos com intervalo de
confiança; classificação em baixo, intermediário e alto risco com base
em thresholds calibrados por DCA; importância de variáveis por SHAP. \\
\textbf{Critérios:} C-index $\geq$ 0,85; time-dependent AUC em 5 anos
$\geq$ 0,88; Brier score integrado $\leq$ 0,15; degradação em validação
externa $\leq$ 5\%; calibração com intercepto $\leq$ 0,10 e inclinação
$\geq$ 0,90. \\
\textbf{Confiança:} 0,83 -- Estudo multicêntrico robusto com validação
externa temporal e análise SHAP abrangente, limitado pela natureza
retrospectiva, exclusão de prontuários incompletos e ausência de dados
genômicos/radiômicos que poderiam elevar adicionalmente a acurácia.

% --- 10. REFERÊNCIA ABNT ---
\subsection{Referência ABNT}
\begin{verbatim}
TANAKA, Hiroshi et al. Machine learning for prediction of oral cancer
survival. Oral Surgery, Oral Medicine, Oral Pathology and Oral
Radiology, New York, v. 138, n. 2, p. 201-215, 2024. DOI:
10.1016/j.oooo.2024.01.010.
\end{verbatim}
\endinput
